A large language model (Claude, Anthropic) was used as a research and engineering assistant for this project, under
the supervision of the authors: to extract and transcribe the proof of Proposition~5.9 from the book, to implement
and test the numerical library and the experiment scripts, to search and verify the bibliography (every reference
was checked against the arXiv, Crossref or Semantic Scholar records), to draft the theory notes (each note was
reviewed by independent adversarial passes checking the constants and the consistency with the cited sources,
and then corrected), and to format this report. All mathematical statements were re-derived and all numerical
claims come from the scripts in the repository (except the cross-check with an independent implementation of
Appendix~\ref{app:repro}, which uses private code), which are deterministic given the seeds reported in
Appendix~\ref{app:repro}, the version of the Newton solver that computes the logistic minimizers and the number of
BLAS threads (with the solver of the cached runs and the default multithreaded BLAS, the fifteen logistic problems of
Section~\ref{sec:exp-B3}, which include those of Sections~\ref{sec:exp-constants}--\ref{sec:exp-A2}, and their
minimizers are rebuilt bit for bit; the corrected solver changes three of these minimizers, and single-threaded BLAS
the data in their last bits), except that a run of the $\kappa$ sweeps of Section~\ref{sec:exp-A4} that exceeds
its CPU-time budget (\nolinkurl{VRLAB_MAX_TIME}, Appendix~\ref{app:repro}) is dropped from the fits, so that which points
enter could depend on the machine and on the budget; in the reported sweeps every run reached its target within the
budget, except the runs of L-Katyusha without negative momentum, which diverge. Representative prompts:
\emph{``Transcribe the proof of Proposition 5.9 as a chain of labelled inequalities''}, \emph{``Try to refute this note:
check every constant''}, \emph{``Verify that this reference exists and that its main result is stated correctly''}.
